\section{Introduction}

Network protocol standards are a key foundation for interoperability and
security across independent implementations. Protocols such as
TLS~\cite{rfc5246,rfc8446}, MQTT~\cite{mqtt311,mqtt5}, and
QUIC~\cite{rfc9000,rfc9001,rfc9002} specify not only how endpoints parse
messages, but also how they negotiate parameters, update state, authenticate
peers, and handle exceptional inputs. A deviation from these requirements does
not necessarily become an exploitable vulnerability immediately, but it can
still introduce practical risk. For example, an implementation may accept a
field value forbidden by the standard, omit a state-dependent check, incorrectly
reject a valid message, or trigger abnormal error-handling behavior under a
specific configuration. Such deviations may cause interoperability failures and
may further enable downgrade attacks, denial-of-service conditions, or other
security risks.

Detecting these deviations is difficult because standards and code usually
express the same behavior in different forms. Protocol standards are
natural-language documents whose requirements are distributed across tables,
field descriptions, state-machine descriptions, compatibility notes, and
exception clauses. RFC-style keywords can indicate requirement
strength~\cite{rfc2119,rfc8174}, and prior work has used RFC-derived rules,
parser specifications, and large-language-model agents to analyze protocol
implementations~\cite{chen2022ribdetector,zheng2025parcleanse,zheng2025rfcaudit,song2026protocolguard}.
However, these approaches usually analyze rules, interfaces, or protocol
behaviors as their main objects, and their granularity is relatively coarse.
They therefore struggle to precisely capture the correspondence between
individual fields, their related conditions, and the implementation behavior
that realizes them.

% Existing techniques only partially solve this problem. Formal validation can 
% provide strong assurance through modeling standards and implementation, but 
% it requires a costly protocol model.  Fuzzing and
% conformance tests exercise real implementations, but their oracles are often
% not linked to individual standard requirements. Static analysis can inspect
% code, but usually lacks standard semantics. LLMs can help connect prose and
% code, but unconstrained LLM judgments are not reliable conformance evidence.
% Effective checking therefore needs a bounded obligation, relevant code context,
% and an explicit way to report insufficient evidence.

This paper focuses on field-level protocol requirements. Many actionable
requirements specify that a field may appear only in a particular state, that
its value must fall within an allowed range, that it must match another field,
that it depends on negotiated parameters, or that it must trigger rejection
under an invalid combination. We represent such requirements as \ruleform{}:
under condition $C$, field $f$ must satisfy action or constraint $A$. This form
is not a complete formal semantics, but it provides a clear boundary for
checking standard text against implementation code.

We present \tool{}, a conformance-checking framework for natural-language
protocol standards and source-code implementations. \tool{} first extracts the
above rules, then retrieves syntax-aware code contexts related to field
parsing, value validation, state dependencies, and error handling, and finally
uses iterative reasoning to determine whether the implementation satisfies the
requirement. Each result is reported as \textsc{Consistent},
\textsc{Inconsistent}, or \textsc{Insufficient}, with supporting or missing
evidence recorded explicitly. For suspicious inconsistent results, \tool{}
further attempts to generate executable validation tests.

We evaluate \tool{} on real-world implementations of TLS/DTLS, MQTT, QUIC, and
CoAP. \tool{} produced 204 submitted conformance-deviation reports; maintainers
confirmed or fixed 126, rejected 9 as false positives, and 69 remain pending.
Two confirmed reports were assigned previously unknown CVEs. Under GPT-5.4
pricing, the average model cost is \$1.122 per submitted report.

Our main contributions are:

\begin{itemize}[leftmargin=*]
  \item We formulate protocol conformance checking as a field-centric alignment
  problem between standards and implementations, and design \tool{}. The
  framework extracts structured field rules from natural-language standards and
  aligns them with implementation code contexts to identify conformance
  deviations in implementations.

  \item We introduce an auditable reporting workflow that turns model-assisted
  analysis into maintainer-oriented defect reports and reproducible validation
  artifacts, reducing reliance on raw LLM verdicts and supporting manual
  confirmation.

  \item We evaluate \tool{} on real-world implementations of TLS/DTLS, MQTT,
  QUIC, and CoAP, producing 204 submitted reports, 126 maintainer-confirmed or
  fixed reports, 9 rejected false positives, 69 pending reports, and two
  previously unknown CVEs.
\end{itemize}
